\phantomsection\addcontentsline{toc}{section}{Cybersecurity}
\ECchapter{16}{Cybersecurity}{How can engineers protect connected systems?}{ECBlue}

\ECObjectives{
\begin{itemize}
\item Define confidentiality, integrity and availability and relate them to physical systems.
\item Explain a defence-in-depth architecture for a connected engineering system.
\item Build and justify a simple cyber-risk analysis from threats, vulnerabilities and consequences.
\end{itemize}}

\begin{multicols}{2}
\begin{engineeringnews}
Industrial systems, vehicles, hospitals and energy networks increasingly depend on connected
software. Cybersecurity is therefore a safety and business-continuity issue as well as an
information-technology issue: a digital attack can produce physical consequences.
\end{engineeringnews}

\begin{ecarticle}
A connected engineering system contains sensors, controllers, software, communication links and
actuators. Each interface can become an entry point for accidental failure or malicious action.
Cybersecurity aims to preserve three central properties: confidentiality prevents unauthorised
disclosure, integrity prevents unauthorised modification, and availability keeps the service usable.

Risk depends on more than the existence of a threat. Engineers identify valuable assets, possible
attackers, vulnerabilities, likelihood and consequences. A weak password on a non-critical display
is not the same risk as an unprotected remote command to a water-treatment pump. Safety analysis
and cyber-risk analysis must therefore be coordinated.

No single protection is sufficient. Defence in depth combines several independent layers. Devices
should authenticate users and other devices. Communications can be encrypted, software should be
updated, and networks should be segmented so that an office computer cannot directly control a
critical machine. Access rights follow the principle of least privilege: each account receives only
the permissions required for its task.

Monitoring is also necessary. Logs, alarms and anomaly detection can reveal unexpected traffic or
commands. When an incident occurs, operators need a prepared response: isolate affected equipment,
maintain a safe physical state, restore trusted software and verify operation. Offline backups help
recovery, but they must be tested rather than merely stored.

Security creates engineering trade-offs. Encryption and authentication use processing time and
energy. Updates may interrupt production, while legacy equipment may not support modern protocols.
A secure design therefore begins with requirements and threat modelling rather than with a protective
tool added at the end. Human factors matter too: clear procedures and training reduce errors and
social-engineering attacks.
\end{ecarticle}

\begin{tcolorbox}[title=\sffamily\bfseries Engineering Insight,colback=yellow!10,colframe=ECOrange]
Safety and cybersecurity can conflict. An emergency door should open quickly for evacuation, yet
access control should prevent intrusion. The engineering task is to define priority rules, degraded
modes and evidence that both requirements are satisfied.
\end{tcolorbox}

\begin{engineeringfocus}
\textbf{Defence-in-depth architecture}

\begin{center}
\resizebox{\linewidth}{!}{%
\begin{tikzpicture}[font=\scriptsize,>=Latex,node distance=3.5mm]
\node[draw=ECBlue,fill=ECLightBlue,rounded corners,align=center] (user) {authorised\\user};
\node[draw=ECGreen,fill=ECGreen!10,rounded corners,align=center,right=of user] (auth) {identity and\\authentication};
\node[draw=ECPurple,fill=ECPurple!10,rounded corners,align=center,right=of auth] (gate) {firewall and\\segmentation};
\node[draw=ECOrange,fill=ECOrange!10,rounded corners,align=center,below=of gate] (ctrl) {industrial\\controller};
\node[draw=ECRed,fill=ECRed!8,rounded corners,align=center,left=of ctrl] (plant) {physical process\\safe state};
\node[draw=ECBlue,fill=ECLightBlue,rounded corners,align=center,left=of plant] (mon) {logging and\\monitoring};
\draw[->,thick] (user)--(auth); \draw[->,thick] (auth)--(gate); \draw[->,thick] (gate)--(ctrl);
\draw[->,thick] (ctrl)--(plant); \draw[->,thick,dashed] (plant)--(mon); \draw[->,thick,dashed] (mon)-| (auth.south);
\end{tikzpicture}%
}
\end{center}

If one layer fails, another should still reduce the probability or consequence of compromise.
\end{engineeringfocus}

\begin{keyfigures}
\begin{tabularx}{\linewidth}{@{}Xr@{}}
Core security objectives & 3\\
Preferred access principle & least privilege\\
Critical architecture measure & segmentation\\
Recovery requirement & tested backup\\
Essential lifecycle activity & updates
\end{tabularx}
\end{keyfigures}

\begin{tcolorbox}[title=\sffamily\bfseries Effect of layered safeguards,colback=white,colframe=ECBlue]
\begin{center}
\begin{tikzpicture}
\begin{axis}[ybar,width=.94\linewidth,height=51mm,ylabel={Residual risk index},ymin=0,ymax=110,
symbolic x coords={No protection,Strong passwords,Updates,Network segmentation,Monitoring,Backup and recovery},xtick=data,
x tick label style={rotate=42,anchor=east,font=\scriptsize},yticklabel style={font=\scriptsize},grid=major]
\addplot table[x=measure,y=risk,col sep=comma]{data/EC16/risk_reduction.csv};
\end{axis}
\end{tikzpicture}
\end{center}
\footnotesize Illustrative risk reduction: safeguards are cumulative, but risk never becomes exactly
zero and each layer must be maintained and periodically verified.
\end{tcolorbox}

\begin{tcolorbox}[title=\sffamily\bfseries Risk reasoning,colback=white,colframe=ECRed]
\begin{center}
\resizebox{\linewidth}{!}{%
\begin{tikzpicture}[font=\scriptsize,>=Latex,node distance=5mm]
\node[draw,rounded corners,fill=ECRed!8] (t) {threat};
\node[draw,rounded corners,fill=ECOrange!10,right=of t] (v) {vulnerability};
\node[draw,rounded corners,fill=ECPurple!10,right=of v] (e) {unwanted event};
\node[draw,rounded corners,fill=ECBlue!10,right=of e] (c) {consequence};
\draw[->,thick] (t)--(v); \draw[->,thick] (v)--(e); \draw[->,thick] (e)--(c);
\end{tikzpicture}%
}
\end{center}
Safeguards may reduce likelihood, limit consequences, improve detection or shorten recovery time.
\end{tcolorbox}

\begin{vocabulary}
\begin{tabularx}{\linewidth}{@{}lX@{}}
confidentiality & confidentialité\\
integrity & intégrité\\
availability & disponibilité\\
threat & menace\\
vulnerability & vulnérabilité\\
authentication & authentification\\
encryption & chiffrement\\
network segmentation & segmentation du réseau\\
least privilege & moindre privilège\\
backup & sauvegarde
\end{tabularx}
\end{vocabulary}

\begin{applicationexercise}
The graph assigns a residual-risk index of 100 to an unprotected system and 11 after all listed
safeguards have been introduced.
\begin{enumerate}
\item Calculate the percentage reduction in residual risk.
\item Which safeguard produces the largest absolute reduction between two consecutive bars?
\item Explain why the final index is not zero.
\item Rank confidentiality, integrity and availability for a connected greenhouse and justify your order.
\item Propose one verification test for the local fail-safe mode during a network outage.
\end{enumerate}
\end{applicationexercise}

\begin{questions}
\item Define confidentiality, integrity and availability.
\item Why should cyber risk be linked to physical consequences?
\item Give four layers of a defence-in-depth strategy.
\item What is the purpose of network segmentation?
\item Why must backups and incident procedures be tested?
\item Distinguish a preventive safeguard from a recovery safeguard.
\end{questions}

\begin{engineeringchallenge}
Secure a connected greenhouse containing temperature sensors, an irrigation pump and remote access.
Identify assets, two threats and two vulnerabilities. Build a small risk matrix, then propose
authentication, segmentation, monitoring and a safe local mode. Finish with a verification plan and a
justified recommendation to the operator.
\end{engineeringchallenge}

\begin{speaking}
Is a fully connected system always better than an isolated one? Compare functionality, maintenance,
security and resilience using one engineering example and one quantified criterion.
\end{speaking}
\end{multicols}

\subsection*{Sources}
\ECsource{NIST -- Cybersecurity Framework}{https://www.nist.gov/cyberframework}
\ECsource{European Union Agency for Cybersecurity}{https://www.enisa.europa.eu/}